\chapter{2013年江西卷（文）}

\section{单选题}

\begin{problem}
复数 \(z = \i(-2 - \i)\)（\(\i\) 为虚数单位）在复平面内所对应的点在（\quad）

\choices
    {第一象限}
    {第二象限}
    {第三象限}
    {第四象限}
\end{problem}

\begin{answer}
D
\end{answer}

\begin{solution}
\(z=\i(-2-\i)=-2\i-\i^{2}=1-2\i\)，所以其对应点的横坐标为正，纵坐标为负，位于第四象限，故选 D.
\end{solution}

\begin{problem}
若集合 \(A = \{x \in \R \mid ax^2 + ax + 1 = 0\}\) 中只有一个元素，则 \(a =\)（\quad）

\choices
    {4}
    {2}
    {0}
    {0 或 4}
\end{problem}

\begin{answer}
A
\end{answer}

\begin{solution}
当 \(a=0\) 时，方程化为 \(1=0\)，无解，不符合题意；当 \(a\ne0\) 时，方程为关于 \(x\) 的一元二次方程，其判别式为
\[
\Delta=a^{2}-4a=a(a-4)
\]
集合中只有一个元素等价于方程有一个实根，因此 \(\Delta=0\). 结合 \(a\ne0\)，得 \(a=4\)，故选 A.
\end{solution}

\begin{problem}
若 \(\sin \frac{\alpha}{2} = \frac{\sqrt{3}}{3}\)，则 \(\cos \alpha =\)（\quad）

\choices
    {\(-\frac{2}{3}\)}
    {\(-\frac{1}{3}\)}
    {\( \frac{1}{3} \)}
    {\( \frac{2}{3} \)}
\end{problem}

\begin{answer}
C
\end{answer}

\begin{solution}
由半角公式，
\[
\cos\alpha=1-2\sin^{2}\frac{\alpha}{2}
 =1-2\left(\frac{\sqrt{3}}{3}\right)^{2}
 =1-\frac{2}{3}=\frac{1}{3}
\]
故选 C.
\end{solution}

\begin{problem}
集合 \(A = \{2, 3\}\)，\(B = \{1, 2, 3\}\)，从 \(A\)，\(B\) 中各任意取一个数，则这两数之和等于 4 的概率是（\quad）

\choices
    {\( \frac{2}{3} \)}
    {\( \frac{1}{2} \)}
    {\( \frac{1}{3} \)}
    {\( \frac{1}{6} \)}
\end{problem}

\begin{answer}
C
\end{answer}

\begin{solution}
从 \(A\) 中取数、从 \(B\) 中取数共有 \(2\times3=6\) 个等可能结果. 其中和为 \(4\) 的结果为 \((2,2)\) 和 \((3,1)\)，共有 \(2\) 个，因此所求概率为
\[
\frac{2}{6}=\frac{1}{3}
\]
故选 C.
\end{solution}

\begin{problem}
总体由编号为 \(01\)，\(02\)，…，\(19\)，\(20\) 的 20 个个体组成. 利用下面的随机数表选取 5 个个体，选取方法从随机数表第 \(1\) 行的第 \(5\) 列和第 \(6\) 列数字开始由左到右依次选取两个数字，则选出来的第 \(5\) 个个体的编号为（\quad）

7816 6572 0802 6314 0702 4369 9728 0198 3204 9234 4935 8200 3623 4869 6938 7481

\choices
    {08}
    {07}
    {02}
    {01}
\end{problem}

\begin{answer}
D
\end{answer}

\begin{solution}
从第 1 行第 5、6 列开始，每两个数字组成一个编号，依次得到
\(65\)，\(72\)，\(08\)，\(02\)，\(63\)，\(14\)，\(07\)，\(02\)，\(43\)，\(69\)，\(97\)，\(28\)，\(01\)，\(98\)，\(32\)，\(04\)，\(\ldots\)
其中只有 \(01\) 至 \(20\) 的编号有效，并且重复编号不再选取. 因此依次选出的个体编号为
\(08\)，\(02\)，\(14\)，\(07\)，\(01\).
第 5 个个体的编号为 \(01\)，故选 D.
\end{solution}

\begin{problem}
下列选项中，使不等式 \(x < \frac{1}{x} < x^2\) 成立的 \(x\) 的取值范围是（\quad）

\choices
    {\((- \infty, -1)\)}
    {\((-1,0)\)}
    {\((0,1)\)}
    {\((1, +\infty)\)}
\end{problem}

\begin{answer}
A
\end{answer}

\begin{solution}
由题意知 \(x\ne0\). 分别解两个不等式：
\[
x<\frac{1}{x}
\iff \frac{x^{2}-1}{x}<0
\iff x\in(-\infty,-1)\cup(0,1)
\]
以及
\[
\frac{1}{x}<x^{2}
\iff \frac{1-x^{3}}{x}<0
\iff x\in(-\infty,0)\cup(1,+\infty)
\]
两解集的交集为 \((-\infty,-1)\)，故选 A.
\end{solution}

\begin{problem}
阅读如图所示程序框图，如果输出 \(i = 4\)，那么在空白的判断框中应填入的条件是（\quad）

\texfigure[float=inline,max width=0.34\linewidth]{img_repaint/2013/jiangxi_liberal/q7_fig1.tex}
\FloatBarrier

\choices
    {\(S < 8\)}
    {\(S < 9\)}
    {\(S < 10\)}
    {\(S < 11\)}
\end{problem}

\begin{answer}
B
\end{answer}

\begin{solution}
由程序框图逐步计算. 初始时 \(i=1\)，\(S=0\)，第一次循环后 \(i=2\)，此时按偶数分支得 \(S=2i+1=5\)；第二次循环后 \(i=3\)，按奇数分支得 \(S=2i+2=8\)；第三次循环后 \(i=4\)，按偶数分支得 \(S=2i+1=9\).

要使输出为 \(i=4\)，应在 \(S=8\) 时继续循环而在 \(S=9\) 时停止，所以判断条件为 \(S<9\)，故选 B.
\end{solution}

\begin{problem}
一几何体的三视图如图所示，则该几何体的体积为


\bitmapfigure[float=inline,max width=0.7\linewidth]{img/2013/jiangxi_liberal/q8_fig1.png}
\FloatBarrier

（\quad）

\choices
    {\(200 + 9\pi\)}
    {\(200 + 18\pi\)}
    {\(140 + 9\pi\)}
    {\(140 + 18\pi\)}
\end{problem}

\begin{answer}
A
\end{answer}

\begin{solution}
由三视图可知，该几何体由一个长、宽、高分别为 \(10\)，\(4\)，\(5\) 的长方体和一个直径为 \(6\)、高为 \(2\) 的半圆柱组成. 长方体体积为
\[
10\times4\times5=200
\]
半圆柱的底面半径为 \(3\)，体积为
\[
\frac{1}{2}\pi\times3^{2}\times2=9\pi
\]
所以几何体体积为 \(200+9\pi\)，故选 A.
\end{solution}

\begin{problem}
已知点 \(A(2,0)\)，抛物线 \(C: x^2 = 4y\) 的焦点为 \(F\)，射线 \(FA\) 与抛物线 \(C\) 相交于点 \(M\)，与其准线相交于点 \(N\)，则 \(|FM|: |MN| =\)（\quad）

\choices
    {\(2:\sqrt{5}\)}
    {\(1:2\)}
    {\( 1 : \sqrt{5} \)}
    {\(1:3\)}
\end{problem}

\begin{answer}
C
\end{answer}

\begin{solution}
抛物线 \(x^{2}=4y\) 的焦点为 \(F(0,1)\)，准线为 \(y=-1\). 设射线 \(FA\) 上点的参数方程为
\((x,y)=(2t,1-t)\)，\(t\geqslant0\).
点 \(A(2,0)\) 对应 \(t=1\). 代入抛物线方程得
\[
4t^{2}=4(1-t)\iff t^{2}+t-1=0
\]
由于 \(M\) 在射线 \(FA\) 上，取正根
\[
t=\frac{\sqrt5-1}{2}
\]
准线 \(y=-1\) 对应 \(1-t=-1\)，即 \(t=2\). 因为同一直线上距离与参数差成正比，
\[
\frac{|FM|}{|MN|}=\frac{t}{2-t}=\frac{1}{\sqrt5}
\]
故 \(|FM|:|MN|=1:\sqrt5\)，故选 C.
\end{solution}

\begin{problem}
如图，已知 \(l_{1} \perp l_{2}\)，圆心在 \(l_{1}\) 上，半径为 \(1 \mathrm{~m}\) 的圆 \(O\) 在 \(t = 0\) 时与 \(l_{2}\) 相切于点 \(A\)，圆 \(O\) 沿 \(l_{1}\) 以 \(1 \mathrm{~m} / \mathrm{s}\) 的速度匀速向上移动，圆被直线 \(l_{2}\) 所截上方圆弧长记为 \(x\)，令 \(y = \cos x\)，则 \(y\) 与时间 \(t\)（\(0 \leqslant t \leqslant 1\)，单位：\(\mathrm{s}\)）的函数 \(y = f(t)\) 的图象大致为（\quad）

\bitmapfigure[float=inline,max width=0.65\linewidth]{img/2013/jiangxi_liberal/q10_fig1.png}
\FloatBarrier

\choices
    {\choicebitmap[width=0.15\paperwidth]{img/2013/jiangxi_liberal/q10_opt_a.png}}
    {\choicebitmap[width=0.15\paperwidth]{img/2013/jiangxi_liberal/q10_opt_b.png}}
    {\choicebitmap[width=0.15\paperwidth]{img/2013/jiangxi_liberal/q10_opt_c.png}}
    {\choicebitmap[width=0.15\paperwidth]{img/2013/jiangxi_liberal/q10_opt_d.png}}
\end{problem}

\begin{answer}
B
\end{answer}

\begin{solution}
取 \(l_{1}\) 为竖直方向. 圆心从与 \(l_{2}\) 相切的位置向上移动 \(t\) 后，圆心到 \(l_{2}\) 的距离为 \(1-t\). 设圆弧所对的圆心角为 \(x\)，则由直角三角形得
\[
\frac{x}{2}=\arccos(1-t)
\]
因此
\[
y=\cos x=\cos\bigl(2\arccos(1-t)\bigr)
 =2(1-t)^{2}-1=2t^{2}-4t+1
\]
在 \(0\leqslant t\leqslant1\) 上，函数从 \(1\) 减到 \(-1\)，且
\(y'=4t-4\leqslant0\)，\(y''=4>0\).
所以图象单调下降且向上凸，与选项 B 相符，故选 B.
\end{solution}

\section{填空题}

\begin{problem}
若曲线 \(y = x^{\alpha} + 1\)（\(\alpha \in \R\)）在点 \((1,2)\) 处的切线经过坐标原点，则 \( \alpha = \)\fillinblank{}.
\end{problem}

\begin{answer}
2
\end{answer}

\begin{solution}
曲线在 \((1,2)\) 处的切线斜率为
\[
y'\big|_{x=1}=\alpha
\]
该切线还经过原点，因此其斜率为
\[
\frac{2-0}{1-0}=2
\]
故 \(\alpha=2\).
\end{solution}

\begin{problem}
某住宅小区计划植树不少于 100 棵，若第一天植 2 棵，以后每天植树的棵数是前一天的 2 倍，则需要的最少天数 \(n\)（\(n \in \N^{*}\)）等于\fillinblank{}.
\end{problem}

\begin{answer}
6
\end{answer}

\begin{solution}
第 \(j\) 天植树 \(2^{j}\) 棵，前 \(n\) 天植树总数为等比数列前 \(n\) 项和
\[
2+2^{2}+\cdots+2^{n}=2^{n+1}-2
\]
要求 \(2^{n+1}-2\geqslant100\). 由于
\(2^{5+1}-2=62<100\)，\(2^{6+1}-2=126\geqslant100\)
所以最少需要 \(n=6\) 天.
\end{solution}

\begin{problem}
设 \( f(x) = \sqrt{3} \sin 3x + \cos 3x \)，若对任意实数 \( x \) 都有 \( |f(x)| \leqslant a \)，则实数 \( a \) 的取值范围是\fillinblank{}.
\end{problem}

\begin{answer}
\(a \in [2,+\infty)\)
\end{answer}

\begin{solution}
由柯西不等式，
\[
|f(x)|=|\sqrt3\sin3x+\cos3x|
 \leqslant\sqrt{(\sqrt3)^{2}+1^{2}}\sqrt{\sin^{2}3x+\cos^{2}3x}=2
\]
当 \(3x=\frac{\pi}{3}\) 时，\(\sin3x=\frac{\sqrt3}{2}\)，\(\cos3x=\frac12\)，从而 \(f(x)=2\)，故最大值确为 \(2\). 因此对任意实数 \(x\) 都有 \(|f(x)|\leqslant a\) 当且仅当 \(a\geqslant2\)，即
\[
a\in[2,+\infty)
\]
\end{solution}

\begin{problem}
若圆 \(C\) 经过坐标原点和点 \((4,0)\)，且与直线 \(y = 1\) 相切，则圆 \(C\) 的方程是\fillinblank{}.
\end{problem}

\begin{answer}
\((x-2)^{2} + \left(y + \frac{3}{2}\right)^{2} = \frac{25}{4}\)
\end{answer}

\begin{solution}
设圆心为 \((h,k)\)，半径为 \(r\). 由于圆经过 \((0,0)\) 和 \((4,0)\)，有
\(h^{2}+k^{2}=r^{2}\)，\((h-4)^{2}+k^{2}=r^{2}\).
两式相减得 \(h=2\)，所以圆心为 \((2,k)\)，且 \(r^{2}=4+k^{2}\). 圆与直线 \(y=1\) 相切，故圆心到该直线的距离等于半径：
\[
|k-1|=\sqrt{4+k^{2}}
\]
平方得
\[
(k-1)^{2}=4+k^{2}\iff k=-\frac32
\]
于是 \(r^{2}=4+\frac94=\frac{25}{4}\)，圆的方程为
\[
\left(x-2\right)^{2}+\left(y+\frac32\right)^{2}=\frac{25}{4}
\]
\end{solution}

\begin{problem}
如图，正方体的底面与正四面体的底面在同一平面 \(\alpha\) 上，且 \(AB \parallel CD\)，则直线 \(EF\) 与正方体的六个面所在的平面相交的平面个数为\fillinblank{}.
\FigureLayoutDeclare{img/2013/jiangxi_liberal/q15_fig1.png}{11.0cm}{35bp 29bp 6bp 8bp}
\bitmapfigure[max width=0.66\linewidth]{img/2013/jiangxi_liberal/q15_fig1.png}
\end{problem}

\begin{answer}
4
\end{answer}

\begin{solution}
设正四面体顶点 \(F\) 在平面 \(\alpha\) 上的射影为正三角形 \(CDE\) 的中心 \(G\). 则直线 \(EF\) 在 \(\alpha\) 上的射影为 \(EG\)，而在正三角形中，中线 \(EG\perp CD\). 因此
\[
EF\perp CD
\]
又因 \(AB\parallel CD\)，所以 \(EF\perp AB\). 于是直线 \(EF\) 与正方体的两个底面棱平行于 \(AB\) 的侧面所在平面相交，与两个底面棱垂直于 \(AB\) 的侧面所在平面平行，不相交.

此外，\(EF\) 在底面平面 \(\alpha\) 上经过点 \(E\)，并且不平行于正方体的上底面所在平面，故还与上下底面所在的两个平面相交. 所以相交的平面共有 \(2+2=4\) 个.
\end{solution}

\section{解答题}

\begin{problem}
正项数列 \(\{a_{n}\}\) 满足： \(a_{n}^{2} - (2n - 1)a_{n} - 2n = 0\).
\begin{enumerate}
\item 求数列 \( \{a_{n}\} \) 的通项公式；
\item 令 \( b_{n} = \frac{1}{(n + 1)a_{n}} \)，求数列 \( \{b_{n}\} \) 的前 \( n \) 项和 \( T_{n} \).
\end{enumerate}
\end{problem}

\begin{solution}
\begin{enumerate}
\item 原方程关于 \(a_n\) 的判别式为
\[
\Delta=(2n-1)^{2}+8n=(2n+1)^{2}
\]
因此
\[
a_n=\frac{2n-1\pm(2n+1)}{2}
\]
即 \(a_n=2n\) 或 \(a_n=-1\). 由于数列为正项数列，取正根，得
\[
a_n=2n
\]
\item 由 \(a_n=2n\)，
\[
b_n=\frac{1}{(n+1)a_n}=\frac{1}{2n(n+1)}
 =\frac12\left(\frac1n-\frac1{n+1}\right)
\]
所以
\[
T_n=\sum_{k=1}^{n}b_k
 =\frac12\sum_{k=1}^{n}\left(\frac1k-\frac1{k+1}\right)
 =\frac12\left(1-\frac1{n+1}\right)
 =\frac{n}{2(n+1)}
\]
\end{enumerate}
\end{solution}

\begin{problem}
在 \(\triangle ABC\) 中，角 \(A\)，\(B\)，\(C\) 的对边分别为 \(a\)，\(b\)，\(c\). 已知 \(\sin A \sin B + \sin B \sin C + \cos 2B = 1\).
\begin{enumerate}
\item 求证： \(a\)，\(b\)，\(c\) 成等差数列；
\item 若 \(C = \frac{2\pi}{3}\)，求 \(\frac{a}{b}\) 的值.
\end{enumerate}
\end{problem}

\begin{solution}
\begin{enumerate}
\item 由已知条件，
\[
\sin B(\sin A+\sin C)=1-\cos2B=2\sin^{2}B
\]
三角形内角满足 \(0<B<\pi\)，所以 \(\sin B>0\)，从而
\[
\sin A+\sin C=2\sin B
\]
由正弦定理，\(\sin A:\sin B:\sin C=a:b:c\)，故
\[
a+c=2b
\]
即 \(a\)，\(b\)，\(c\) 成等差数列.
\item 当 \(C=\frac{2\pi}{3}\) 时，\(A+B=\frac\pi3\). 令 \(t=\tan\frac A2\)，则 \(0<A<\frac\pi3\)，故 \(t>0\). 由 \(\sin A+\sin C=2\sin B\) 得
\[
2\sin A-\sqrt3\cos A+\frac{\sqrt3}{2}=0
\]
代入
\(\sin A=\frac{2t}{1+t^{2}}\)，\(\cos A=\frac{1-t^{2}}{1+t^{2}}\)
并整理得
\[
3\sqrt3t^{2}+8t-\sqrt3=0
\]
其根为 \(t=\frac{\sqrt3}{9}\) 和负根；由 \(t>0\)，取 \(t=\frac{\sqrt3}{9}\). 因而
\(\sin A=\frac{3\sqrt3}{14}\)，\(\sin B=\sin\left(\frac\pi3-A\right)=\frac{5\sqrt3}{14}\).
再由正弦定理，
\[
\frac ab=\frac{\sin A}{\sin B}=\frac35
\]
\end{enumerate}
\end{solution}

\begin{problem}
小波以游戏方式决定是去打球，唱歌还是去下棋，游戏规则为：以 \(O\) 为起点，再从 \(A_{1}\)，\(A_{2}\)，\(A_{3}\)，\(A_{4}\)，\(A_{5}\)，\(A_{6}\)（如图）这 \(6\) 个点中任取两点分别为终点得到两个向量，记这两个向量的数量积为 \(X\)，若 \(X > 0\) 就去打球，若 \(X = 0\) 就去唱歌，若 \(X < 0\) 就去下棋.
\begin{enumerate}
\item 写出数量积 \(X\) 的所有可能取值；
\item 分别求小波去下棋的概率和不去唱歌的概率.
\end{enumerate}
\bitmapfigure[max width=0.34\linewidth]{img/2013/jiangxi_liberal/q18_fig1.png}
\end{problem}

\begin{solution}
\begin{enumerate}
\item 由图可取
\(A_1=(1,0)\)，\(A_2=(1,-1)\)，\(A_3=(0,-1)\)，\(A_4=(0,1)\)
\(A_5=(-1,1)\)，\(A_6=(-1,0)\).
任取两个不同点，计算相应向量的数量积. 当点对为 \((A_2,A_5)\) 时，\(X=-2\)，共 1 种；当点对为 \((A_1,A_5)\)，\((A_1,A_6)\)，\((A_2,A_4)\)，\((A_2,A_6)\)，\((A_3,A_4)\)，\((A_3,A_5)\) 时，\(X=-1\)，共 6 种；当点对为 \((A_1,A_3)\)，\((A_1,A_4)\)，\((A_3,A_6)\)，\((A_4,A_6)\) 时，\(X=0\)，共 4 种；当点对为 \((A_1,A_2)\)，\((A_2,A_3)\)，\((A_4,A_5)\)，\((A_5,A_6)\) 时，\(X=1\)，共 4 种.
所以 \(X\) 的所有可能取值为 \(-2\)，\(-1\)，\(0\)，\(1\).
\item 两点的无序选法共有 \(\mathrm C_6^2=15\) 种；若按有序选取，则每个无序点对对应两个等可能的有序结果，概率相同. 由上述计数，\(X<0\) 有 \(1+6=7\) 种，故去下棋的概率为
\[
P(X<0)=\frac7{15}
\]
不去唱歌等价于 \(X\ne0\)，因此
\[
P(X\ne0)=1-\frac4{15}=\frac{11}{15}
\]
\end{enumerate}
\end{solution}

\begin{problem}
如图，直四棱柱 \(ABCD - A_{1}B_{1}C_{1}D_{1}\) 中，\(AB \parallel CD\)，\(AD \perp AB\)，\(AB = 2\)，\(AD = \sqrt{2}\)，\(AA_{1} = 3\)，\(E\) 为 \(CD\) 上一点，\(DE = 1\)，\(EC = 3\).
\begin{enumerate}
\item 证明： \(BE \perp\) 平面 \(BB_{1}C_{1}C\)；
\item 求点 \(B_{1}\) 到平面 \(EA_{1}C_{1}\) 的距离.
\end{enumerate}
\bitmapfigure[max width=0.46\linewidth]{img/2013/jiangxi_liberal/q19_fig1.png}
\end{problem}

\begin{solution}
\begin{enumerate}
\item 建立空间直角坐标系，令
\(A=(0,0,0)\)，\(B=(2,0,0)\)，\(D=(0,\sqrt2,0)\)，\(C=(4,\sqrt2,0)\)
\(A_1=(0,0,3)\)，\(B_1=(2,0,3)\)，\(C_1=(4,\sqrt2,3)\)，\(E=(1,\sqrt2,0)\).
则
\(\overrightarrow{BE}=(-1,\sqrt2,0)\)，\(\overrightarrow{BC}=(2,\sqrt2,0)\)，\(\overrightarrow{BB_1}=(0,0,3)\).
有
\(\overrightarrow{BE}\cdot\overrightarrow{BC}=-2+2=0\)，\(\overrightarrow{BE}\cdot\overrightarrow{BB_1}=0\).
直线 \(BC\) 与 \(BB_1\) 是平面 \(BB_1C_1C\) 内相交的两条直线，因此 \(BE\perp\) 平面 \(BB_1C_1C\).
\item 取
\(\overrightarrow{EA_1}=(-1,-\sqrt2,3)\)，\(\overrightarrow{EC_1}=(3,0,3)\).
平面 \(EA_1C_1\) 的一个法向量可取
\[
(-\sqrt2,4,\sqrt2)
\]
因为
\[
\overrightarrow{EA_1}\times\overrightarrow{EC_1}=3(-\sqrt2,4,\sqrt2)
\]
所以其方程为
\[
-\sqrt2(x-1)+4(y-\sqrt2)+\sqrt2z=0
\]
即
\[
-\sqrt2x+4y+\sqrt2z-3\sqrt2=0
\]
点 \(B_1(2,0,3)\) 到此平面的距离为
\[
\frac{|-2\sqrt2+3\sqrt2-3\sqrt2|}{\sqrt{2+16+2}}
 =\frac{2\sqrt2}{2\sqrt5}=\frac{\sqrt{10}}5
\]
\end{enumerate}
\end{solution}

\begin{problem}
椭圆 \(C\)：\(\frac{x^2}{a^2} + \frac{y^2}{b^2} = 1\)（\(a > b > 0\)）的离心率 \(e = \frac{\sqrt{3}}{2}\)，\(a + b = 3\).
\begin{enumerate}
\item 求椭圆 \(C\) 的方程；
\item 如图所示，\(A\)，\(B\)，\(D\) 是椭圆 \(C\) 的顶点，\(P\) 是椭圆 \(C\) 上除顶点外的任意一点，直线 \(DP\) 交 \(x\) 轴于点 \(N\)，直线 \(AD\) 交 \(BP\) 于点 \(M\)，设 \(BP\) 的斜率为 \(k\)，\(MN\) 的斜率为 \(m\)，证明：\(2m-k\) 为定值.
\end{enumerate}
\bitmapfigure[max width=0.42\linewidth]{img/2013/jiangxi_liberal/q20_fig1.png}
\end{problem}

\begin{solution}
\begin{enumerate}
\item 由离心率，
\(e^{2}=\frac{c^{2}}{a^{2}}=\frac34\)，\(b^{2}=a^{2}-c^{2}=\frac14a^{2}\).
因 \(a\)，\(b>0\)，有 \(b=\frac a2\). 再由 \(a+b=3\) 得 \(a=2\)，\(b=1\)，故椭圆方程为
\[
\frac{x^{2}}4+y^{2}=1
\]
\item 由图取椭圆顶点 \(A=(-2,0)\)，\(B=(2,0)\)，\(D=(0,1)\). 设直线 \(BP\) 的斜率为 \(k\). 因 \(P\) 不是顶点，\(k\) 为有限非零数，且 \(2k-1\ne0\)、\(2k+1\ne0\)：\(k=0\) 时 \(P=A\)，\(k=\frac12\) 时 \(P\) 为下顶点，\(k=-\frac12\) 时 \(P=D\)，均与题意矛盾. 因此下列除法均有意义. 直线 \(BP\) 的方程为 \(y=k(x-2)\)，与椭圆联立，除去交点 \(B\) 后得
\[
P=\left(\frac{2(4k^{2}-1)}{4k^{2}+1},-\frac{4k}{4k^{2}+1}\right)
\]
直线 \(AD\) 的方程为 \(y=\frac x2+1\)，所以
\[
M=AD\cap BP=\left(\frac{2(2k+1)}{2k-1},\frac{4k}{2k-1}\right)
\]
直线 \(DP\) 与 \(x\) 轴交于
\[
N=\left(\frac{2(2k-1)}{2k+1},0\right)
\]
于是 \(MN\) 的斜率为
\[
m=\frac{\dfrac{4k}{2k-1}}{\dfrac{2(2k+1)}{2k-1}-\dfrac{2(2k-1)}{2k+1}}
 =\frac{2k+1}{4}
\]
因此
\[
2m-k=2\cdot\frac{2k+1}{4}-k=\frac12
\]
为定值.
\end{enumerate}
\end{solution}

\begin{problem}
设函数 \( f(x) = \begin{cases}
\frac{1}{a} x, & 0 \leqslant x \leqslant a,\\
\frac{1}{1 - a} (1 - x), & a < x \leqslant 1,
\end{cases} \) \( a \) 为常数且 \( a \in (0,1) \).
\begin{enumerate}
\item 当 \(a = \frac{1}{2}\) 时，求 \(f\left(f\left(\frac{1}{3}\right)\right)\)；
\item 若 \(x_0\) 满足 \(f(f(x_0)) = x_0\)，但 \(f(x_0) \neq x_0\)，则称 \(x_0\) 为 \(f(x)\) 的二阶周期点，证明：函数 \(f(x)\) 有且仅有两个二阶周期点，并求出二阶周期点 \(x_1\)，\(x_2\)；
\item 对于（2）中的 \(x_{1}\)，\(x_{2}\)，设 \(A\left(x_{1}, f(f(x_{1}))\right)\)，\(B\left(x_{2}, f(f(x_{2}))\right)\)，\(C\left(a^{2}, 0\right)\)，记 \(\triangle ABC\) 的面积为 \(S(a)\)，求 \(S(a)\) 在区间 \(\left[\frac{1}{3}, \frac{1}{2}\right]\) 上的最大值和最小值.
\end{enumerate}
\end{problem}

\begin{solution}
\begin{enumerate}
\item 当 \(a=\frac12\) 时，\(\frac13\leqslant a\)，故
\[
f\left(\frac13\right)=2\cdot\frac13=\frac23
\]
又 \(\frac23>a\)，所以
\[
f\left(f\left(\frac13\right)\right)=f\left(\frac23\right)=2\left(1-\frac23\right)=\frac23
\]
\item 记 \(d=1+a-a^{2}\). 分段计算复合函数：
\[
f(f(x))=
\begin{cases}
\dfrac{x}{a^{2}},&0\leqslant x\leqslant a^{2},\\[4pt]
\dfrac{a-x}{a(1-a)},&a^{2}<x\leqslant a,\\[4pt]
\dfrac{x-a}{(1-a)^{2}},&a<x<1-a+a^{2},\\[4pt]
\dfrac{1-x}{a(1-a)},&1-a+a^{2}\leqslant x\leqslant1.
\end{cases}
\]
在第一段令 \(f(f(x))=x\)，只有固定点 \(x=0\)；第二段给出
\[
x=\frac{a}{1+a-a^{2}}=\frac ad;
\]
第三段给出固定点
\[
x=\frac{1}{2-a};
\]
第四段给出
\[
x=\frac{1}{1+a-a^{2}}=\frac1d
\]
由于 \(0<a<1\)，有
\(a^{2}<\frac ad<a\)，\(a<\frac1{2-a}<1-a+a^{2}\)，\(1-a+a^{2}<\frac1d<1\)
所以上述三个非零解分别确实落在对应分段内. 其中 \(x=0\) 和 \(x=\frac1{2-a}\) 满足 \(f(x)=x\)，不属于二阶周期点；而
\(x_1=\frac ad\)，\(x_2=\frac1d\)
分别落在第二、四段，并满足
\(f(x_1)=x_2\)，\(f(x_2)=x_1\).
由于 \(0<a<1\)，有 \(x_1\ne x_2\)，故它们正好是两个二阶周期点，且函数没有其他二阶周期点.
\item 对上述两点，\(f(f(x_i))=x_i\)，所以
\(A=(x_1,x_1)\)，\(B=(x_2,x_2)\)，\(C=(a^{2},0)\).
由三角形面积的行列式公式，且 \(x_2>x_1\)，
\[
S(a)=\frac12\left|(x_2-x_1)(-x_1)-(x_2-x_1)(a^{2}-x_1)\right|
 =\frac12(x_2-x_1)a^{2}
\]
又
\[
x_2-x_1=\frac{1-a}{d}
\]
故
\[
S(a)=\frac{a^{2}(1-a)}{2(1+a-a^{2})}
\]
求导得
\[
S'(a)=\frac{a(2-2a-2a^{2}+a^{3})}{2(1+a-a^{2})^{2}}
\]
当 \(a\in[\frac13,\frac12]\) 时，分母为正，且
\[
2-2a-2a^{2}+a^{3}\geqslant2-1-\frac12=\frac12>0
\]
所以 \(S(a)\) 在该区间上单调递增. 因此
\(S_{\min}=S\left(\frac13\right)=\frac1{33}\)，\(S_{\max}=S\left(\frac12\right)=\frac1{20}\).
\end{enumerate}
\end{solution}
